% ═══════════════════════════════════════════════════════════════
% Section 10: Conclusion
% ═══════════════════════════════════════════════════════════════

\section{Conclusion}
\label{sec:conclusion}

Brain-computer interfaces are moving from laboratory prototypes to commercial medical
devices. Security frameworks designed for information technology are necessary for these
devices but not sufficient: a compromised neural device can act on a person, can expose
data that cannot be reissued, and can outlast the party that supports it. That is the case
for neurosecurity. Because the harms include violated consent and effects that cannot be
undone, and because claims about brains are easily overstated, a neurosecurity framework
also needs neuroethics, both to decide what counts as harm and to limit what the framework
itself may claim.

This paper presented the \qif framework, a proposed and unvalidated system comprising an
addressing model for what a device touches, a threat catalogue (TARA), an impact-only CVSS
v4.0 extension (NISS, six metrics, with three scoring states), and the Neural Impact Chain,
a threat-modeling method that associates techniques with diagnostic category references.
In the \SnapshotTechniques{}-technique snapshot analysed here, CVSS cannot express the
primary impact of 59 techniques and misses part of it for 30 more; 71.6\% of techniques have
a recorded therapeutic analog, showing how often the mechanism of an attack is the mechanism
of a therapy; and 56.0\% are in a direct or indirect association class.

This revision also corrected the paper. It withdrew the hourglass stack and its chokepoint,
retired an evidence status that conflated unrelated kinds of evidence, replaced case studies
that did not match the registry, and reported which published figures could not be
reproduced. The framework's limits remain: no empirical validation on BCI hardware, no
clinical review of the category associations, no expert review of the addressing model, and
single-author scoring. These are invitations to collaborate, not reasons to delay
publication.

The framework, catalogue, NISS specification and scoring data are released under the
Apache 2.0 license. We invite review and correction from the neuroscience, neuroethics,
cybersecurity and clinical communities; a finding that an entry is wrong is as useful as a
finding that it holds.
